\documentclass[10pt,twocolumn]{article}
\usepackage[T1]{fontenc}
\usepackage[utf8]{inputenc}
\usepackage{lmodern}
\usepackage[margin=1.8cm,top=2.4cm,bottom=2cm,columnsep=0.6cm]{geometry}
\usepackage{fancyhdr}
\usepackage{titlesec}
\usepackage{booktabs}
\usepackage{amsmath,amssymb,amsfonts}
\usepackage{microtype}
\usepackage{xcolor}
\usepackage{mdframed}
\usepackage{parskip}
\usepackage{enumitem}
\usepackage{hyperref}

\definecolor{rulered}{RGB}{180,30,30}
\definecolor{boxgray}{RGB}{245,245,245}
\definecolor{keywordgray}{RGB}{100,100,100}

\pagestyle{fancy}
\fancyhf{}
\fancyhead[L]{\textsc{\small Claude Mag}}
\fancyhead[C]{\small\textit{Machine Learning Research Digest}}
\fancyhead[R]{\small 2026-08-24}
\fancyfoot[C]{\small\thepage}
\renewcommand{\headrulewidth}{0.8pt}

\titleformat{\section}{\normalsize\bfseries\color{rulered}}{}{0pt}{}%
  [\vspace{-4pt}\color{rulered}\rule{\columnwidth}{0.4pt}]
\titlespacing{\section}{0pt}{8pt}{4pt}

\hypersetup{colorlinks=true,urlcolor=rulered,linkcolor=black}

\begin{document}
\setlength{\parindent}{0pt}
\setlength{\parskip}{4pt}

{\small\color{keywordgray}\textbf{Keywords:}
  vision-language models \textbullet{}
  benchmark \textbullet{}
  visual decision making \textbullet{}
  constraint satisfaction}\par\vspace{4pt}

{\LARGE\bfseries\raggedright
  ComplexityWorld: Benchmarking Vision-Language
  Models on Verifiable Visual Decision
  Making}\par\vspace{4pt}

{\large\itshape\raggedright
  390 Tasks Across 39 Visual Worlds Show That Most VLMs Cannot
  Extract Constraints from Images and Satisfy Them Jointly ---
  Only One Model Exceeds 75\%}\par\vspace{6pt}

{\small\textbf{By} ComplexityWorld Authors}\par\vspace{2pt}

{\small\color{keywordgray}arXiv:2608.07584 \textbullet{} Preprint, August 2026}
\par\vspace{2pt}\noindent{\color{rulered}\rule{\columnwidth}{0.6pt}}\vspace{6pt}

Many real-world AI tasks require more than perceiving what an image
contains: a model must use visual evidence to make a decision whose
components jointly satisfy global constraints. ComplexityWorld
introduces a benchmark of 390 tasks across 39 domain-inspired visual
worlds and 29 decision categories, each scored by an executable
verifier that accepts any feasible solution. Under direct visual
inference, all evaluated VLMs except GPT-5.6-Sol fall below 40\%
verifier acceptance rate (VAR), while GPT-5.6-Sol reaches 75.6\%.
A 23-point gap between visual and structured-text conditions reveals
that the dominant failure mode is visual constraint extraction, not
reasoning.

\begin{mdframed}[backgroundcolor=boxgray,linecolor=rulered,linewidth=1pt,
    innerleftmargin=6pt,innerrightmargin=6pt,
    innertopmargin=6pt,innerbottommargin=6pt]
\textbf{\small AT A GLANCE}\par\vspace{4pt}
\begin{description}[leftmargin=0pt,labelsep=4pt,itemsep=2pt,parsep=0pt,
    font=\small\bfseries]
  \item[Problem:] VLM benchmarks test perception and QA accuracy but
    not the ability to make constraint-satisfying decisions from visual
    evidence, leaving a large deployment-relevant gap unmeasured.
  \item[Why it matters:] Real-world tasks (scheduling, logistics,
    design) require visually-grounded constraint satisfaction; current
    evaluations do not capture whether models can do this.
  \item[Method:] Generate 390 tasks from hidden structured
    specifications rendered as visual scenes; use executable verifiers
    rather than LLM judges to score feasibility.
  \item[Result:] Only GPT-5.6-Sol exceeds 40\% VAR (reaching 75.6\%);
    all other models remain below 40\%.
  \item[Evidence:] Same tasks in structured-text form achieve $+23$ pp
    average VAR, isolating visual extraction as the key failure mode.
\end{description}
\end{mdframed}

\section{The Big Picture}

Vision-language models have advanced rapidly on perception tasks ---
they can accurately describe image contents, count objects, and answer
factual questions about visual scenes. This progress has driven
optimism about deploying VLMs in real-world decision support: assisting
with medical imaging, factory floor monitoring, logistics planning,
and scientific data analysis.

These real-world tasks share a property that standard benchmarks do
not test: they require making a \emph{decision} from visual evidence,
where the decision must satisfy multi-part constraints. A scheduling
assistant reviewing a visual calendar cannot merely describe meetings;
it must identify conflicts, dependencies, and resource constraints
from the visual layout and produce a feasible schedule satisfying
all of them simultaneously.

This class of problem --- \emph{visual constraint-satisfying decision
making} --- is distinct from visual perception and from verbal
reasoning, and it is largely unmeasured by current evaluation suites.
ComplexityWorld is the first systematic benchmark for this capability.

\section{Why Existing Approaches Fall Short}

\textbf{Perception benchmarks} (VQA, object counting, scene
description) test whether models correctly identify what is present
in an image. They do not test whether models can extract structured
decision-relevant information from visual representations of
constraints.

\textbf{Verbal reasoning benchmarks} (GSM8K, logical reasoning
datasets) test constraint satisfaction in text, but cannot assess
whether visual extraction precedes and enables correct reasoning.

\textbf{Prior VLM decision-making benchmarks} have used LLM-as-judge
evaluation, which is susceptible to hallucination and cannot
distinguish a feasible from a plausible-sounding infeasible solution.
ComplexityWorld's verifier-based evaluation avoids these issues.

\section{Core Mechanism}

\textbf{Task generation pipeline:} Each task starts from a
\emph{hidden structured specification} $\mathcal{S}$ that defines:
\begin{itemize}[itemsep=1pt,parsep=0pt]
  \item A set of objects, resources, or agents with attributes
  \item A set of hard constraints over decision variables
  \item A decision space (assignments, routes, orderings, partitions)
\end{itemize}
The specification is rendered into a visual scene $I$ using a
domain-appropriate renderer (grid-based for scheduling, graph-based
for routing, CAD-based for spatial arrangement). An executable
verifier $\mathcal{V}$ checks whether any response $r$ satisfies
all constraints in $\mathcal{S}$.

\textbf{Key design choices:}
The verifier accepts \emph{any} feasible solution, not just the
canonical one, avoiding benchmark contamination and penalizing models
solely for constraint violations rather than stylistic deviation.

\section{Technical Formulation}

Let $\mathcal{S} = (X, C, \mathcal{D})$ be a task specification with
decision variables $X = \{x_1, \ldots, x_n\}$, constraints
$C = \{c_1, \ldots, c_m\}$ over $X$, and domain $\mathcal{D}$.

A response $r$ parsed from the model's text output is a partial or
complete assignment $\hat{X} = \{\hat{x}_1, \ldots, \hat{x}_n\}$.
Verifier acceptance:
\[
  \text{Accept}(r, \mathcal{S}) = \mathbf{1}\!\left[
    \bigwedge_{j=1}^m c_j(\hat{X}) = \text{TRUE}\right]
\]
The verifier acceptance rate (VAR) for model $M$ over task set $\mathcal{T}$:
\[
  \text{VAR}(M) = \frac{1}{|\mathcal{T}|}\sum_{i \in \mathcal{T}}
    \text{Accept}(M(I_i), \mathcal{S}_i)
\]
The visual constraint extraction gap is measured by the difference
between structured-text and visual-scene performance:
\[
  \Delta_\text{VCE}(M) = \text{VAR}(M, \text{structured})
    - \text{VAR}(M, \text{visual})
\]

\section{Learning or Inference Procedure}

ComplexityWorld is an evaluation benchmark, not a training procedure.
Models are evaluated zero-shot with a standardized prompt that
specifies the task domain, output format, and response constraints.
No few-shot examples are provided in the main evaluation; ablations
with $k=1,3,5$ in-context examples show improvement but do not
close the gap with structured-text performance.

\section{What the Guarantee Says}

The verifier is formally specified for each domain using executable
constraint logic. Verifier correctness is validated by human experts
who confirm that feasibility checks match domain semantics.
This guarantees that VAR measures genuine constraint satisfaction,
not LLM-judge agreement with plausible-sounding responses.

\section{Experimental Findings}

\begin{itemize}[itemsep=2pt,parsep=0pt]
  \item \textbf{Direct visual inference VAR:} GPT-5.6-Sol
    reaches 75.6\%; all other models remain below 40\%.
  \item \textbf{Visual constraint extraction gap:} Average
    $\Delta_\text{VCE} = +23.1$ pp when the same task is
    given in structured text vs.\ as a visual scene.
  \item \textbf{Presentation sensitivity:} Changing visual style
    (color scheme, font, layout) while keeping constraints identical
    shifts VAR by up to 18 pp for the same model, revealing
    brittleness in visual feature extraction.
  \item \textbf{Category variation:} Routing and scheduling domains
    are hardest (mean VAR $<25\%$); assignment and partitioning
    domains are easier (mean VAR $\sim$38\%).
  \item \textbf{Open vs.\ proprietary models:} Best open-weight
    model achieves 31.4\% VAR; best proprietary model (GPT-5.6-Sol)
    achieves 75.6\%, a gap of 44.2 pp.
\end{itemize}

\section{Ablations and Interpretation}

\textbf{In-context examples:} Adding $k=5$ in-context (image, decision)
examples improves VAR by 6.2 pp on average, but narrows the
visual--structured gap by only 1.3 pp, confirming that the bottleneck
is visual extraction, not few-shot learning of the decision format.

\textbf{Chain-of-thought prompting:} Requiring models to reason
step-by-step before producing a decision improves VAR by 4.8 pp
on average, with larger gains on routing ($+8.1$ pp) and scheduling
($+7.3$ pp) tasks that have natural sequential decision structure.

\textbf{Constraint complexity:} VAR decreases sharply with the
number of constraints $m$: from 52\% at $m=2$ to 21\% at $m=6$
for a representative model, suggesting exponential constraint
combinatorics as the dominant difficulty driver.

\textbf{Error analysis:} Manual inspection of 200 failures reveals
three failure modes: misreading a constraint value from the image
(43\%), correctly extracting constraints but producing an infeasible
assignment (38\%), and producing a feasible partial assignment that
misses some variables (19\%).

\section*{Reference}

ComplexityWorld Authors.
\textbf{ComplexityWorld: Benchmarking Vision-Language Models
on Verifiable Visual Decision Making.}
arXiv:2608.07584, August 2026.
\url{https://arxiv.org/abs/2608.07584}

\end{document}
